\section{Introduction}
\label{sec:introduction}

A durable workflow engine runs long-lived application logic and keeps its
state across process restarts. Engines differ in the programming model they
give the application. Dex's model is deliberately small. An application
defines a \emph{Flow}. A Flow is a set of \emph{Step} types. A Step has an
optional \kw{WaitFor} method, which returns what the Step waits for, and a
required \kw{Execute} method, which returns a \emph{StepDecision}: go to
other Steps, end this branch, or close the Flow. Steps read and write
\emph{Attributes}, publish to and consume from \emph{Channels}, and wait on
\emph{Timers}. There is no thread that blocks in the middle of a function,
no lock that is held across a wait, and no shared memory that is read and
written in one atomic step.

A natural question follows: what can this model \emph{not} express? The
usual way to answer is by example. The Workflow Patterns
\cite{vanderaalst2003patterns,russell2006revised} are a catalog of 43
control-flow behaviors, and an engine is rated by how directly it supports
each one \cite{wohed2009patterns}. Dex has such an evaluation: 19 of the 43
control-flow patterns have direct support, 24 have indirect support, and
none is unsupported; each rating comes with a runnable example and an
integration test against a real server.%
\footnote{The evaluation is on the Dex documentation site, under
\emph{Concepts, Correctness and expressiveness, Workflow Patterns
evaluation}.}
A pattern catalog measures breadth against the running system. It does not
say that some \emph{class} of programs is expressible, and it does not say
what ``expressible'' means when the engine interleaves the pieces of a
program in ways the source program never wrote down.

This paper gives the other kind of evidence: two precise claims, each with
a machine-checked proof.

\begin{quote}
\emph{Every well-formed program of a small concurrent imperative language
\Lang{} translates into a Dex Flow that is trace-equivalent to it, assuming
no failures.}
\end{quote}

\begin{quote}
\emph{In the same fragment and under explicit hypotheses, erasing the
repeated external effects from a finite trace with crashes, takeovers,
retries, zombie Workers, and duplicate deliveries gives a failure-free
trace, and every failure-free trace remains possible. For fair runs the
erasure keeps the termination, deadlock, or divergence marker.}
\end{quote}

The first claim is about expressiveness, and \S\ref{sec:language} to
\S\ref{sec:engineering} develop it. The second is about durability, the
reason a workflow engine exists, and \S\ref{sec:failures} develops it.

\paragraph{What ``expressive'' means here.}
\Lang{} is a conventional shared-memory concurrent language: threads with
local variables, shared variables, procedures with recursion, \kw{fork} and
\kw{join}, mutexes, condition variables, \kw{sleep}, a timed \kw{select} on
an external channel, cooperative cancellation, exceptions, external effects,
and program exit. We define a translation $\tr{\cdot}$ from \Lang{} programs
to Dex Flow definitions, a set of observable events on each side, and a
projection $\pi$ from Dex events to \Lang{} events. The claim is two
equalities of trace sets, for every well-formed \Lang{} program $S$:
\begin{align}
\pi\bigl(\FinTr_{\Eadm}(\tr{S})\bigr) &= \FinTr(S)
  && \text{(safety)} \label{eq:safety}\\
\forall w \in \Eadm.\quad
\pi\bigl(\MaxTr^{\mathit{fair}}_{w}(\tr{S})\bigr) &= \MaxTr^{\mathit{fair}}_{w}(S)
  && \text{(liveness)} \label{eq:liveness}
\end{align}
\FinTr{} is the set of finite observable traces. $\MaxTr^{\mathit{fair}}_w$
is the set of traces of fair infinite runs for the input word $w$, each
with a marker: terminated (\mterm), deadlocked (\mdead), or divergent
(\mdiv). \Eadm{} is the admissible environment: it starts the program once
and then sends plain data on the program's external channels.
Equation~\eqref{eq:safety} says that the Flow does nothing the program could
not do, and the program does nothing the Flow could not do. Equation
\eqref{eq:liveness} adds that the two terminate, deadlock, and diverge under
the same conditions, which finite traces alone cannot tell.

So ``Dex is expressive'' means: the Dex core can simulate a language that
has the standard synchronization vocabulary of concurrent programming. The
correspondence uses finite sequences of Dex transitions, some observable
and some silent: an action of the program corresponds to a short sequence
of Dex transitions, and some Dex transitions correspond to no action of
the program. The observable traces are preserved. The translation is
the constructive content. It shows \emph{how} a mutex, a condition variable,
or a cancellable timed receive is built from Steps, Channels, Attributes,
and Timers.

\paragraph{Contributions.}
\begin{enumerate}
\item \textbf{A formal semantics of the Dex core}
  (\S\ref{sec:dex}). It is a labeled transition system written against the
  implementation of the Dex server, with the rules citing the code that
  implements them. (One rule outside the theorem's fragment follows a fix
  that was not yet merged at the pinned commit, and says so.) Three
  adversarial reviews checked it against the code.
\item \textbf{The language \Lang{} and its semantics}
  (\S\ref{sec:language}), including the choices that make an exact
  equivalence possible: a timing model, a well-formedness scope, and
  unforgeable handles.
\item \textbf{The translation} from \Lang{} to Dex
  (\S\ref{sec:translation}): one Step per program point, in
  continuation-passing style.
\item \textbf{The theorems} (\S\ref{sec:theorems}) and \textbf{their proof
  in Lean~4} (\S\ref{sec:proof}): \lean{safety_of_isWellFormed} and
  \lean{liveness_of_isWellFormed}. No proof contains a \texttt{sorry}
  placeholder, the development depends on the Lean core library only, and
  \texttt{\#print axioms} reports
  \lean{propext}, \lean{Classical.choice}, and \lean{Quot.sound} for both
  theorems.
\item \textbf{Non-vacuity} of both theorems: a decidable static lock
  discipline that implies well-formedness, witness programs that together
  use every construct of \Lang{}, and kernel-checked fair maximal traces
  for each of the three markers.
\item \textbf{An account of the proof engineering}
  (\S\ref{sec:engineering}): parallel work packages with fixed statements,
  the gaps that the provers found in the fixed definitions, the fuzz that
  now tests those definitions for inductiveness, and two defects that the
  work exposed in the real server.
\item \textbf{Failure transparency} (\S\ref{sec:failures}): a failure
  model written against the server, with crashes, takeovers, Worker
  failures, zombie Workers, stale results, and StartFlow retries; an
  abstraction map with ghost history under which every failure step is a
  stutter; and the Lean theorems \lean{safety_failure_forward},
  \lean{liveness_failure}, \lean{endToEnd_safety}, and
  \lean{endToEnd_liveness}, with the same three axioms. Kernel-checked
  witnesses apply the safety statement and the forward liveness inclusion
  with the markers \mterm{} and \mdead{}. Counterexamples show that H1
  and H3 are needed and that a re-sent publication is a second request;
  the H4 counterexample also violates H1, of which H4 is a special case in
  the model.
\end{enumerate}

\paragraph{What is not claimed.}
The scope is narrow on purpose, and \S\ref{sec:limitations} discusses each
limit.
\begin{itemize}
\item \emph{Failures, under hypotheses.} The failure theorem covers SYNC
  Step methods only, assumes that no method exhausts its retries and that
  user code ignores the attempt number and similar context, assumes
  finitely many failures for liveness, and keeps the timer order of the
  first theorem as an idealization. Effects may repeat and overlap; the
  theorem does not make them exactly-once.
\item \emph{A model, not the code.} The Dex side of the theorem is a
  hand-written model of the server, validated by review against the source.
  It is not extracted from the implementation.
\item \emph{Well-formed programs only.} A program that can call
  \kw{unlock}, \kw{wait}, \kw{signal}, or \kw{broadcast} without holding the
  mutex is outside the domain, as a program with a data race is outside a
  theorem about data-race-free programs.
\item \emph{A restricted environment.} Inputs are one start and then
  publications to external channels. Remote procedure calls into a Flow,
  stopping a Flow, and the other client operations are specified in the
  semantics but are outside the theorem.
\item \emph{One timing model.} \Lang{} counts a fixed number of clock-stamped
  actions per statement, aligned with Dex. The theorem does not say that a
  program is equivalent to itself under another cost model.
\item \emph{Trace equivalence.} The theorem is not bisimilarity and not full
  abstraction, and it says nothing about performance.
\item \emph{One direction of expressiveness.} We translate \Lang{} into Dex.
  We do not claim that every Dex Flow is an \Lang{} program.
\end{itemize}

\paragraph{Artifacts.}
The semantics is \file{formal/semantics/dex-core-semantics.md} (``the
specification'' below). The first theorem formalizes its revision~8;
revision~9 adds the failure model of Section~11 and leaves every
failure-free rule unchanged. The Lean development is \file{formal/lean/}.
File paths in this paper are relative to the repository root, and Lean
files are under \file{formal/lean/DexExpressiveness/}. Numbers for the
first theorem are for commit \texttt{33e998003} of the development, and
numbers for the failure theorem for commit \texttt{dd449bd57}, which
changes no file of the first proof except for import lines.
